\textbf{Equação diferencial exponencial (eq. 5.10):}
\[
\frac{dN}{dt} = \lambda N.
\]

\textbf{Verificação de que $N(t)=N_0 e^{\lambda t}$ (eq. 5.12) é solução:}

\textbf{Método 1 --- substituição direta:}

Calculando a derivada:
\[
\frac{dN}{dt} = \frac{d}{dt}\left(N_0 e^{\lambda t}\right) = N_0\lambda e^{\lambda t} = \lambda\underbrace{N_0 e^{\lambda t}}_{=N(t)} = \lambda N(t).\quad\checkmark
\]

\textbf{Método 2 --- integração usando $\int du/u = \ln u + C$ (conforme solicitado):}

Separando variáveis na eq. (5.10):
\[
\frac{dN}{N} = \lambda\,dt.
\]

Integrando ambos os lados:
\[
\int\frac{dN}{N} = \int\lambda\,dt \implies \ln N = \lambda t + C,
\]
onde $C$ é a constante de integração.

Aplicando a condição inicial $N(0)=N_0$:
\[
\ln N_0 = \lambda\cdot0 + C \implies C = \ln N_0.
\]

Portanto:
\[
\ln N = \lambda t + \ln N_0 \implies \ln\frac{N}{N_0} = \lambda t \implies \frac{N}{N_0} = e^{\lambda t}.
\]

Concluindo:
\[
\boxed{N(t) = N_0\,e^{\lambda t}.}\quad\blacksquare
\]